
\documentclass{article}
\usepackage{colortbl}
\usepackage{makecell}
\usepackage{multirow}
\usepackage{supertabular}

\begin{document}

\newcounter{utterance}

\centering \large Interaction Transcript for game `codenames', experiment `association\_easy', episode 0 with reka{-}flash{-}3.1{-}t0.0.
\vspace{24pt}

{ \footnotesize  \setcounter{utterance}{1}
\setlength{\tabcolsep}{0pt}
\begin{supertabular}{c@{$\;$}|p{.15\linewidth}@{}p{.15\linewidth}p{.15\linewidth}p{.15\linewidth}p{.15\linewidth}p{.15\linewidth}}
    \# & $\;$A & \multicolumn{4}{c}{Game Master} & $\;\:$B\\
    \hline

    \theutterance \stepcounter{utterance}  
    & & & \multicolumn{2}{p{0.3\linewidth}}{
        \cellcolor[rgb]{0.9,0.9,0.9}{
            \makecell[{{p{\linewidth}}}]{
                \texttt{\tiny{[GM$|$GM]}}
                {'hidden': {'team': ['australia', 'goat', 'watermelon', 'lion', 'korea', 'iran', 'bear', 'plum', 'moose'], 'innocent': ['whisk', 'missile', 'barbie', 'spatula', 'machete', 'blocks', 'yoyo'], 'opponent': ['mosquito', 'hotel', 'flood', 'clouds', 'mansion', 'wasp', 'shack', 'cottage'], 'assassin': ['parakeet']}, 'revealed': {'team': {'team': [], 'innocent': [], 'opponent': [], 'assassin': []}, 'opponent': {'team': [], 'innocent': [], 'opponent': [], 'assassin': []}}}
            }
        }
    }
    & & \\ \\

    \theutterance \stepcounter{utterance}  
    & & \multicolumn{4}{p{0.6\linewidth}}{
        \cellcolor[rgb]{0.9,0.9,0.9}{
            \makecell[{{p{\linewidth}}}]{
                \texttt{\tiny{[P1$\langle$GM]}}
                \texttt{Provide one single{-}word clue that relates to some of your team's words from the list below. You can choose to how many and to which words your clue relates to.} \\
\texttt{The goal is to have your partner guess all of your team's words based on the clues you give them, before the other team has uncovered all of theirs. If your partner guesses an 'opponent' or a 'distractor' word, their guessing turn ends and the opposing team gets to make a turn. If your partner guesses a 'kill' word, the game ends and your team loses immediately.} \\
\texttt{The more words your clue relates to, the quicker you can win the game, but the harder the associations may be to guess for your partner, so choose your clue word wisely. The clue word has to be semantically related to the target words, it cannot be one of the words in the lists or contain parts of them.} \\
\\ 
\texttt{Always give your single{-}word clue and your comma{-}separated list of related target words in the following format and make your answers as short as possible, never include any other text than is required in this form:} \\
\\ 
\texttt{CLUE: <WORD>} \\
\texttt{TARGETS: <TARGETS>} \\
\\ 
\texttt{Your team words are: australia, goat, watermelon, lion, korea, iran, bear, plum, moose.} \\
\texttt{Your opponent's team words are: mosquito, hotel, flood, clouds, mansion, wasp, shack, cottage.} \\
\texttt{Distractor words are: whisk, missile, barbie, spatula, machete, blocks, yoyo.} \\
\texttt{The kill words are: parakeet.} \\
            }
        }
    }
    & \\ \\

    \theutterance \stepcounter{utterance}  
    & \multicolumn{4}{p{0.6\linewidth}}{
        \cellcolor[rgb]{0.9,0.9,0.9}{
            \makecell[{{p{\linewidth}}}]{
                \texttt{\tiny{[P1$\rangle$GM]}}
                \texttt{<reasoning>} \\
\texttt{Let's tackle this clue{-}giving problem. The user's team has words like Australia, goat, watermelon, lion, Korea, Iran, bear, plum, moose. I need to find a single{-}word clue that relates to as many of these as possible without being one of the words themselves or containing parts of them.} \\
\\ 
\texttt{First, let's look for common themes. Australia and lion are animals, but maybe not directly related. Wait, lion is a big cat, and Australia has kangaroos, but maybe not. Let's see: goat, bear, lion, moose—all are animals. Australia is a country, but maybe there's a connection. Wait, maybe "animal" is too broad. Let's check other possibilities.} \\
\\ 
\texttt{Looking at the countries: Australia, Korea, Iran. Maybe something related to geography? Or maybe each country has a national animal? For example, Australia's national animal is the kangaroo, but that's not in the list. Lion is the national animal of some countries, like Senegal. Iran's national animal is the Persian lion. Wait, Iran's word is on the list. So lion is related to Iran. Also, Australia's national animal isn't in the list, but maybe "lion" is connected to Iran. Hmm.} \\
\\ 
\texttt{Alternatively, maybe the clue is about animals. Let's see: goat, bear, lion, moose. Those are all mammals. Australia is a country, but maybe "kangaroo" is a distractor. Wait, the user's team words don't include kangaroo. So maybe "animal" is too vague. Let's see if there's another angle.} \\
\\ 
\texttt{Plum and watermelon are fruits. Wait, plum is a fruit, watermelon is a fruit. But the other words: Australia, goat, lion, Korea, Iran, bear, moose. Maybe "fruit" only relates to plum and watermelon. That's two words. Maybe not enough. Let's check other possibilities.} \\
\\ 
\texttt{Looking at the countries again: Australia, Korea, Iran. Maybe something like "continent"? Australia is a continent, Korea is part of Asia, Iran is in Asia. But that might not be strong enough. Alternatively, maybe "country" but that's too broad.} \\
\\ 
\texttt{Wait, maybe "animal" is the way to go. Let's list the animal words: goat, bear, lion, moose. That's four words. Plus, maybe Australia? Wait,} \\
            }
        }
    }
    & & \\ \\

    \theutterance \stepcounter{utterance}  
    & & & \multicolumn{2}{p{0.3\linewidth}}{
        \cellcolor[rgb]{0.9,0.9,0.9}{
            \makecell[{{p{\linewidth}}}]{
                \texttt{\tiny{[GM$|$GM]}}
                {'player': 'Cluegiver', 'type': 'rambling error', 'utterance': '<reasoning>\nLet\'s tackle this clue-giving problem. The user\'s team has words like Australia, goat, watermelon, lion, Korea, Iran, bear, plum, moose. I need to find a single-word clue that relates to as many of these as possible without being one of the words themselves or containing parts of them.\n\nFirst, let\'s look for common themes. Australia and lion are animals, but maybe not directly related. Wait, lion is a big cat, and Australia has kangaroos, but maybe not. Let\'s see: goat, bear, lion, moose—all are animals. Australia is a country, but maybe there\'s a connection. Wait, maybe "animal" is too broad. Let\'s check other possibilities.\n\nLooking at the countries: Australia, Korea, Iran. Maybe something related to geography? Or maybe each country has a national animal? For example, Australia\'s national animal is the kangaroo, but that\'s not in the list. Lion is the national animal of some countries, like Senegal. Iran\'s national animal is the Persian lion. Wait, Iran\'s word is on the list. So lion is related to Iran. Also, Australia\'s national animal isn\'t in the list, but maybe "lion" is connected to Iran. Hmm.\n\nAlternatively, maybe the clue is about animals. Let\'s see: goat, bear, lion, moose. Those are all mammals. Australia is a country, but maybe "kangaroo" is a distractor. Wait, the user\'s team words don\'t include kangaroo. So maybe "animal" is too vague. Let\'s see if there\'s another angle.\n\nPlum and watermelon are fruits. Wait, plum is a fruit, watermelon is a fruit. But the other words: Australia, goat, lion, Korea, Iran, bear, moose. Maybe "fruit" only relates to plum and watermelon. That\'s two words. Maybe not enough. Let\'s check other possibilities.\n\nLooking at the countries again: Australia, Korea, Iran. Maybe something like "continent"? Australia is a continent, Korea is part of Asia, Iran is in Asia. But that might not be strong enough. Alternatively, maybe "country" but that\'s too broad.\n\nWait, maybe "animal" is the way to go. Let\'s list the animal words: goat, bear, lion, moose. That\'s four words. Plus, maybe Australia? Wait,', 'message': 'Your answer contained more than two lines, please only give one clue and your targets on two separate lines.'}
            }
        }
    }
    & & \\ \\

\end{supertabular}
}

\end{document}
